%! Polynomial Zeros and Factoring — Unit 3, Lesson 3.4 — Guided Notes
% THE in-class document, in the MAIN IDEAS / NOTES shape (LESSON_SHAPE.md §1): the vocabulary
% box, the hook, then ONE two-column guidednotes table. Each row is one idea — a label on the
% left; on the right one or two COMPLETE printed sentences, ONE large pre-drawn display, then a
% two-across grid of numbered problems with real writing room. Problems run 1..18.
% Every graph is pre-drawn and READ — no sketching.
%   I DO   : the teacher reads the definitions, marks up the displays, works problems 1 and 3,
%            then 7 and 11 at the head of their rows.
%   WE DO  : the class works the rest of each grid, every one cold-called.
%   YOU DO : the last row — problems 15–18, alone, for 16 minutes while the teacher circulates.
% The crux is row 4, TWO PATHS: dividing both sides by x loses the zero x = 0 (problems 11, 12,
% 14); problem 13 is the second trap, stopping early at (x^2 - 4)(x^2 + 4). You Do 16 transfers
% the divide trap, 17 the stopping-early trap, 18 both.
%
% PAGE LOCKSTEP with notes_key/: the vocabulary write lines -> \ansline, the hook's write
% lines -> \ansline, and the answer argument of every \pcell, \writespace and \labelbox filled.
% Nothing else differs. Inside a cell `\\` ENDS THE ROW — break lines with \par. A row cannot
% break across pages: the figure and EACH grid row are separate sub-rows (`\\` then `& ...`).
\documentclass[10pt]{article}
\usepackage{precalculus-article}
\usepackage{precalculus-boxes}
\newcommand{\TallMath}[1]{$\displaystyle #1\rule[-1.4em]{0pt}{3.2em}$}

\begin{document}

\pageheader{Unit 3, Lesson 3.4}{Guided Notes}

\vspace{0.12in}

\boxguard
\begin{vocabbox}
Fill in each term as we use it together.
\par\vspace{2pt}
\noindent\textbf{\textcolor{plum}{Zero Product Property:}}\\[1pt]
\writeline\\[3pt]
\noindent\textbf{\textcolor{plum}{Greatest common factor (GCF):}}\\[1pt]
\writeline\\[3pt]
\noindent\textbf{\textcolor{plum}{Factor by grouping:}}\\[1pt]
\writeline\\[3pt]
\noindent\textbf{\textcolor{plum}{Sum / difference of cubes:}}\\[1pt]
\writeline\\[3pt]
\noindent\textbf{\textcolor{plum}{Quadratic form:}}\\[1pt]
\writeline\\[3pt]
\noindent\textbf{\textcolor{plum}{Factored completely:}}\\[1pt]
\writeline
\end{vocabbox}

\vspace{0.1in}

\boxguard
\begin{hookbox}
Jordan solves $x^3 = 4x$ by dividing both sides by $x$: $x^2 = 4$, so $x = 2$ or $x = -2$.
Riley says, ``You lost one.'' The equation has degree 3.

How many solutions should it have? Can you find the one Jordan lost?
\writelines{2}
\end{hookbox}

\small
\begin{guidednotes}
% ── ROW 1: zero product property and the GCF — factor, then set each factor to zero ──────
\mainidea[A product is zero]{Factor First} &
\textbf{Zero Product Property:} if a product equals zero, at least one of its factors equals
zero. So to solve a polynomial equation, \textbf{move every term to one side} so the other side
is 0, \textbf{factor}, and set \textbf{each factor} equal to zero. Always look for a
\textbf{greatest common factor (GCF)} first --- including a power of $x$.
\\
& \begin{minipage}[c]{0.46\linewidth}\centering
\begin{tikzpicture}[>=stealth]
  \node[font=\large] (p) at (0,0) {$p(x) = x\,(x + 4)\,(x - 3)$};
  \draw[->, thick, plum] (-0.55,-0.3) -- (-1.7,-1.1);
  \draw[->, thick, plum] (0.45,-0.3) -- (0,-1.1);
  \draw[->, thick, plum] (1.55,-0.3) -- (1.7,-1.1);
  \node[font=\small] at (-1.7,-1.35) {$x = 0$};
  \node[font=\small] at (0,-1.35) {$x + 4 = 0$};
  \node[font=\small] at (1.7,-1.35) {$x - 3 = 0$};
\end{tikzpicture}

\vspace{4pt}
{\footnotesize Zeros:}\ \labelbox{0.9cm}{}\ \labelbox{0.9cm}{}\ \labelbox{0.9cm}{}
\end{minipage}\hfill
\begin{minipage}[c]{0.5\linewidth}\centering
\begin{tikzpicture}[x=0.42cm, y=0.075cm, >=stealth]
  \draw[linegray, very thin, xstep=1, ystep=5] (-5,-15) grid (4,25);
  \draw[->, thick] (-5,0) -- (4.5,0) node[right] {\small $x$};
  \draw[->, thick] (0,-15) -- (0,27) node[above] {\small $p(x)$};
  \foreach \x in {-4,-3,-2,-1,1,2,3} \node[below, font=\tiny] at (\x,0) {\x};
  \foreach \y in {-10,10,20} \node[left, font=\tiny] at (0,\y) {\y};
  \draw[very thick, plum, domain=-4.45:3.55, samples=70, smooth, variable=\t]
       plot (\t, {\t*(\t+4)*(\t-3)});
  \foreach \z in {-4,0,3} \fill[goldacc] (\z,0) circle (3pt);
\end{tikzpicture}
\end{minipage}
\\
& \notesprompt{Solve. Factor out the GCF first, then set each factor equal to zero.}
\begin{probgrid}{|Y|Y|}
\pcell{1}{$2x^2 - 10x = 0$}{2.0cm}{} &
\pcell{2}{$6x^3 + 9x^2 = 0$. Which solution has multiplicity 2?}{2.0cm}{} \\ \hline
\end{probgrid}
\\ \hline
% ── ROW 2: the factoring checklist — GCF, count the terms, factor completely, check ──────
\mainidea[A checklist]{Factoring Strategy} &
Factor in this order, every time. \stepnum{1}\ Take out the \textbf{GCF}.
\stepnum{2}\ \textbf{Count the terms}: two terms --- a difference of squares or a sum or
difference of cubes; three terms --- a trinomial, or a \textbf{quadratic form} such as
$x^4 - 5x^2 + 4$, where $u = x^2$; four terms --- \textbf{grouping}.
\stepnum{3}\ \textbf{Factor completely}: test every factor again until none can be factored over
the real numbers. \stepnum{4}\ \textbf{Check} by multiplying back.
\\
& \begin{center}
\renewcommand{\arraystretch}{1.5}
\begin{tabular}{|l|l|}
\hline
\rowcolor{lilac}\textbf{Two-term pattern} & \textbf{Factors as} \\ \hline
Difference of squares & $a^2 - b^2 = (a - b)(a + b)$ \\ \hline
Difference of cubes & $a^3 - b^3 = (a - b)(a^2 + ab + b^2)$ \\ \hline
Sum of cubes & $a^3 + b^3 = (a + b)(a^2 - ab + b^2)$ \\ \hline
Sum of squares & $a^2 + b^2$ \ is \ \labelbox{3.0cm}{} over the reals \\ \hline
\end{tabular}

\vspace{3pt}
{\footnotesize Grouping: $x^3 + 2x^2 - 9x - 18 = x^2(x + 2) - 9(x + 2) = (x + 2)(x^2 - 9)$ ---
and then step 3.}
\end{center}
\\
& \notesprompt{Factor completely. Name the strategy at each step.}
\begin{probgrid}{|Y|Y|}
\pcell{3}{$x^3 + 2x^2 - 9x - 18$}{2.0cm}{} &
\pcell{4}{$2x^3 - 50x$}{2.0cm}{} \\ \hline
\end{probgrid}
\\
& \begin{probgrid}*{|Y|Y|}
\pcell{5}{$27x^3 + 1$}{2.0cm}{} &
\pcell{6}{$x^4 - 5x^2 + 4$ \quad (let $u = x^2$)}{2.0cm}{} \\ \hline
\end{probgrid}
\\ \hline
% ── ROW 3: every factor gives solutions — real or not; degree n gives n ─────────────────
\mainidea[Real or not]{Every Solution} &
Set \textbf{every} factor equal to zero, even one whose graph never reaches the $x$-axis:
$x^2 + 9 = 0$ gives $x = \pm 3i$, and a quadratic factor that will not factor is solved with
the quadratic formula. Non-real solutions come in \textbf{conjugate pairs} $a \pm bi$, and a
polynomial equation of degree $n$ has exactly $n$ solutions, counting multiplicity. Only the
\textbf{real} solutions are $x$-intercepts.
\\
& \begin{minipage}[c]{0.46\linewidth}\centering
\begin{tikzpicture}[x=0.75cm, y=0.12cm, >=stealth]
  \draw[linegray, very thin, xstep=1, ystep=2] (-3,-8) grid (3,10);
  \draw[->, thick] (-3,0) -- (3.5,0) node[right] {\small $x$};
  \draw[->, thick] (0,-8) -- (0,11.5) node[above] {\small $y$};
  \foreach \x in {-2,-1,1,2} \node[below, font=\tiny] at (\x,0) {\x};
  \foreach \y in {-6,-4,-2,2,4,6,8} \node[left, font=\tiny] at (0,\y) {\y};
  \draw[very thick, plum, domain=-2.33:2.33, samples=70, smooth, variable=\t]
       plot (\t, {\t*\t*\t*\t - 3*\t*\t - 4});
  \node[font=\scriptsize, anchor=west, fill=white, inner sep=1pt] at (0.6,10.3) {$y = x^4 - 3x^2 - 4$};
\end{tikzpicture}
\end{minipage}\hfill
\begin{minipage}[c]{0.5\linewidth}
\raggedright
The graph crosses the $x$-axis \labelbox{1.0cm}{} times.

\vspace{4pt}
The equation $x^4 - 3x^2 - 4 = 0$ has degree 4, so it has \labelbox{1.0cm}{} solutions.

\vspace{4pt}
The others are a \labelbox{2.6cm}{} pair.
\end{minipage}
\\
& \notesprompt{Solve. Give every solution, real and non-real.}
\begin{probgrid}{|Y|Y|}
\pcell{7}{$x^4 - 3x^2 - 4 = 0$ \quad (the graph above)}{2.2cm}{} &
\pcell{8}{$x^2 + 9 = 0$. Why does $y = x^2 + 9$ have no $x$-intercepts?}{2.2cm}{} \\ \hline
\end{probgrid}
\\
& \begin{probgrid}*{|Y|Y|}
\pcell{9}{$x^3 + 2x^2 + 4x + 8 = 0$}{2.4cm}{} &
\pcell{10}{$x^3 - 8 = 0$ \quad (use the quadratic formula on the quadratic factor)}{2.4cm}{} \\ \hline
\end{probgrid}
\\ \hline
% ── ROW 4: THE CRUX — divide by x, or move to one side and factor? ──────────────────────
\mainidea[Divide, or factor?]{Two Paths} &
\textbf{Never divide both sides by $x$} (or by any power of $x$). Dividing assumes $x \ne 0$, so
it throws away the solution $x = 0$ without a word. Move every term to one side, factor out the
power of $x$, and keep it as a factor. And \textbf{finish}: a factor such as $x^2 - 4$ is not
done until it is $(x - 2)(x + 2)$.
\\
& \begin{minipage}[c]{0.64\linewidth}\centering
\begin{tikzpicture}[>=stealth]
  \node[draw=redacc, very thick, fill=redbg, rounded corners=3pt, align=left,
        text width=3.0cm, font=\footnotesize] at (0,0)
        {\textbf{Path A: divide by $x$}\\[3pt]
         $x^3 = 9x$\\[2pt]
         $x^2 = 9$\\[2pt]
         $x = \pm 3$\\[2pt]
         \textit{two solutions}};
  \node[draw=plum, very thick, fill=lilac, rounded corners=3pt, align=left,
        text width=3.6cm, font=\footnotesize] at (3.7,0)
        {\textbf{Path B: one side, factor}\\[3pt]
         $x^3 - 9x = 0$\\[2pt]
         $x(x^2 - 9) = 0$\\[2pt]
         $x(x - 3)(x + 3) = 0$\\[2pt]
         $x = 0,\ 3,\ -3$};
\end{tikzpicture}

\vspace{4pt}
{\footnotesize Path A lost $x =$} \labelbox{0.9cm}{}
\end{minipage}\hfill
\begin{minipage}[c]{0.34\linewidth}\centering
\begin{tikzpicture}[x=0.36cm, y=0.12cm, >=stealth]
  \draw[linegray, very thin, xstep=1, ystep=5] (-4,-12) grid (4,12);
  \draw[->, thick] (-4,0) -- (4.5,0) node[right] {\scriptsize $x$};
  \draw[->, thick] (0,-12) -- (0,13.5) node[above] {\scriptsize $y$};
  \foreach \x in {-3,3} \node[below, font=\tiny] at (\x,0) {\x};
  \draw[very thick, plum, domain=-3.45:3.45, samples=60, smooth, variable=\t]
       plot (\t, {\t*\t*\t - 9*\t});
  \foreach \z in {-3,0,3} \fill[goldacc] (\z,0) circle (2.6pt);
  \node[font=\tiny, fill=white, inner sep=1pt] at (2.3,-10.5) {$y = x^3 - 9x$};
\end{tikzpicture}
\end{minipage}
\\
& \notesprompt{Two answers disagree. Decide which is right, and fix the other.}
\begin{probgrid}{|Y|Y|}
\pcell{11}{Which path is right? Substitute $x = 0$ into $x^3 = 9x$. How many solutions should a
degree-3 equation have?}{2.2cm}{} &
\pcell{12}{Solve $2x^4 = 8x^2$ both ways. What does dividing by $2x^2$ lose? Count the solutions
with multiplicity.}{2.2cm}{} \\ \hline
\end{probgrid}
\\
& \begin{probgrid}*{|Y|Y|}
\pcell{13}{Ana writes $x^4 - 16 = (x^2 - 4)(x^2 + 4)$, calls it factored completely, and says
$x^4 - 16 = 0$ has solutions $2$ and $-2$ only. Fix both claims.}{2.4cm}{} &
\pcell{14}{Solve $x^3 + 4x = 4x^2$. Then show what dividing both sides by $x$ would have
lost.}{2.4cm}{} \\ \hline
\end{probgrid}
\\ \hline
% ── LAST ROW: YOU DO — alone, while the teacher circulates (16 min) ─────────────────────
\mainidea[You do, alone]{Factor and Solve} &
Solve each equation by factoring. Move every term to one side, factor completely, and give
every solution --- real and non-real --- with its multiplicity.
\\
& \notesprompt{On your own. Show your work.}
\begin{probgrid}{|Y|Y|}
\pcell{15}{$3x^3 - 12x = 0$}{2.4cm}{} &
\pcell{16}{$5x^3 = 20x^2$. How many solutions, counting multiplicity?}{2.4cm}{} \\ \hline
\end{probgrid}
\\
& \begin{probgrid}*{|Y|Y|}
\pcell{17}{Factor completely, then solve: $x^4 - 81 = 0$.}{3.0cm}{} &
\pcell{18}{\begin{tikzpicture}[x=0.55cm, y=0.04cm, baseline=(current bounding box.north)]
  \draw[linegray, very thin, xstep=1, ystep=10] (-3,-30) grid (3,30);
  \draw[->, thick] (-3,0) -- (3.4,0) node[right] {\scriptsize $x$};
  \draw[->, thick] (0,-30) -- (0,33) node[above] {\scriptsize $p(x)$};
  \foreach \x in {-2,-1,1,2} \node[below, font=\tiny] at (\x,0) {\x};
  \draw[very thick, plum, domain=-2.32:2.32, samples=70, smooth, variable=\t]
       plot (\t, {\t*\t*\t*\t*\t - \t*\t*\t - 12*\t});
\end{tikzpicture}\par
The graph of $p(x) = x^5 - x^3 - 12x$ shows three $x$-intercepts. Find all five zeros of $p$.
Where are the other two?}{3.0cm}{} \\ \hline
\end{probgrid}
\\ \hline
\end{guidednotes}

% ── THE NOTES END AT THE YOU DO ROW; AP Practice follows; homework is DeltaMath ──────

\end{document}
